\startmode[single]

\setvariables[meta]
  [title={Hammingcode},
   subtitle={Zusammenfassung},
   date={2024-10-18},
   author={Niklas von Hirschfeld},
  ]
\setupinteraction[title={\getvariable{meta}{title}}]

\starttext

\startfrontmatter
\component c_titlepage
% \component c_toc
\stopfrontmatter

\setuppagenumber[number=1]

\startbodymatter
\stopmode

\definecolor[blue][r=0.2, g=0.5, b=0.8]
\definecolor[green][r=0.2, g=0.7, b=0.3]
\definecolor[yellow][r=0.8, g=0.7, b=0.2]
\definecolor[orange][r=1, g=0.5, b=0.2]

% Content here:

\startchapter[title={(7, 4) Hamming-Code}]

\startframedtext[location=middle,align=middle,width=\dimexpr\textwidth+\rightmarginwidth+\rightmargindistance\relax,frame=off]
\startnarrower[left,right]
Der {\em (7, 4) Hamming-Code} ist eine möglichkeit Bit-Fehler bei der
Übertragung (z.B. in einem Netzwerk) zu erkennen. Dies ist notwendig, da die
Datenübertragung in der Regel nicht fehlerfrei ist.
\stopnarrower
\stopframedtext


Der Hamming-Code nutzt dafür ein ähnliches Verfahren wie bei einem
Parität\sidenote{Ein Bit, welches angibt ob die anzahl der einsen in dem Byte
gerade oder ungerade war. Sollte sich dies ändern (z.B. durch fehlerhafte
übertragung) wird diese Fehler sichtbar. Man kann aber nicht sagen, was falsch
ist, {\bf nur} das etwas an diesem Byte falsch ist.}. Allerdings wird nicht nur
ein Bit genutzt, sonder ganze drei. Dadurch kommt auch der Name zustande
(\in{s. Abb.}[fig:hamming:text]). Die \color[blue]{$7$} steht für die Anzahl
insgesammt zu übertragene Bits und die \color[green]{$4$} wie viele davon für
Informationen genutzt werden. Also bleiben \color[orange]{$3$} übrig, mit
welchen man {\bf einen} exakten Fehler bestimmen kann. Im gegensatzt zu einem
Paritätsbit kann auch die genaue Position ermittelt werden. Sollte es
allerdings zu mehr als einem Fehler kommen kann die Position nicht mehr
bestimmt werden.

\margintext{

\startplacefigure[title={Zustandekommen der Zahlen 7 und 4},reference={fig:hamming:text},width=\rightmarginwidth,location=here]
\startalignment[middle]
% Hamming code title
{\tfa 
\color[yellow]{(}
\color[blue]{7}
,
\color[green]{4}
\color[yellow]{) Hamming-Code}

\vskip 0.4cm

% Binary sequence in Hamming code
\color[green]{1011}\color[orange]{010}

\vskip 0.4cm

% Mathematical equation for (7 - 4 = 3)
\color[blue]{$7$}$-$
\color[green]{$4$}$=$
\color[orange]{$3$}
}
\stopalignment
\stopplacefigure

}

Damit können insgesammt $57\%$ ($4\div7$) der übertragenen Daten für die
eigentliche Information genutzt werden, der Rest gilt zur Fehlerbestimmung.


\startsectionlevel[title={Berechnung des Hamming-Code}]

\startbuffer[mp:fig:hammingocde:bits:darstellung]
\startMPpage

u := 1cm;
% Farben definieren
color blue; blue := (0.2, 0.3, 0.7);
color green; green := (0.3, 0.6, 0.3);
color orange; orange := (0.8, 0.5, 0.2);
color yellow; yellow := (1, 0.8, 0.3);
color gray; gray := (0.6, 0.6, 0.6);


% Kreise definieren
path cone, ctwo, cthree;
cone := fullcircle scaled (3 * u) shifted (u * dir 0);
ctwo := fullcircle scaled (3 * u) shifted (u * dir 120);
cthree := fullcircle scaled (3 * u) shifted (u * dir 240);

% Kreise mit Transparenz und Farben füllen
fill cone withcolor blue withtransparency (1, 0.6);
fill ctwo withcolor orange withtransparency (1, 0.6);
fill cthree withcolor green withtransparency (1, 0.6);

draw cone withcolor blue;
draw ctwo withcolor orange;
draw cthree withcolor green;

label("$p_1$", (u * 1.5 * dir 0));
label("$p_2$", (u * 1.5 * dir 120));
label("$p_3$", (u * 1.5 * dir 240));

label("$d_1$", (u * dir 60));
label("$d_2$", (u * dir 180));
label("$d_3$", (u * dir 300));
label("$d_4$", origin);


\stopMPpage
\stopbuffer

\margintext{
\placefigure[here][mp:fig:hammingocde:bits:darstellung]{Visualisierung der Aufteilung unter den Bits}{
\typesetbuffer[mp:fig:hammingocde:bits:darstellung]
}}


Die $3$ und $4$ kommen zu stande, da jeder Paritätbit {\bf drei} der anderen
Bits \quotation{Überwacht}. Den drei Paritätbits ($p_1, p_2
\text{ und } p_3$) werden also die vier Bits ($d_1, d_2, d_3 \text{ und }
d_4$) zugewiesen (\in{s. Abb.}[mp:fig:hammingocde:bits:darstellung]).

Um jetzt die Paritätsbits zu berechnen wird die anzahl der Einsen in den
zugehörigen Bits gezählt. Sollte diese gerade sein, ist der Paritätbit eine
Null, sollte die Anzahl ungerade sein ein Eins.

\startbuffer[fig:mp:hammingcode:bits:zuweisen]
\startMPpage

u := 1cm;
% Farben definieren
color blue; blue := (0.2, 0.3, 0.7);
color green; green := (0.3, 0.6, 0.3);
color orange; orange := (0.8, 0.5, 0.2);
color yellow; yellow := (1, 0.8, 0.3);
color gray; gray := (0.6, 0.6, 0.6);


% Kreise definieren
path cone, ctwo, cthree;
cone := fullcircle scaled (3 * u) shifted (u * dir 0);
ctwo := fullcircle scaled (3 * u) shifted (u * dir 120);
cthree := fullcircle scaled (3 * u) shifted (u * dir 240);

% Kreise mit Transparenz und Farben füllen
fill cone withcolor blue withtransparency (1, 0.6);
fill ctwo withcolor orange withtransparency (1, 0.6);
fill cthree withcolor green withtransparency (1, 0.6);

draw cone withcolor blue;
draw ctwo withcolor orange;
draw cthree withcolor green;

label("$p_1$", (u * 1.5 * dir 0));
label("$p_2$", (u * 1.5 * dir 120));
label("$p_3$", (u * 1.5 * dir 240));

label("$1$", (u * dir 60));
label("$0$", (u * dir 180));
label("$1$", (u * dir 300));
label("$1$", origin);

\stopMPpage
\stopbuffer

\margintext{
\placefigure[here][fig:mp:hammingcode:bits:zuweisen]{Verteilung der Bits in der Besipielrechnung}{
\typesetbuffer[fig:mp:hammingcode:bits:zuweisen]
}}

Nun berechnen wir die Paritätbits an dem Beispiel der der Bits: $1011$. Wir
weisen die Bits wie in \in{Abbildung}[fig:mp:hammingcode:bits:zuweisen] dargestellt zu.

Bei $p_1$ zählen wir drei Bits mit $1$, da drei ungerade ist wird dieser
Paritätbit auf $p_1 = 1$ gesetzt. Für $p_2 \text{ und } p_3$ zählen wir jeweils
nur zwei, somit werden sie auf $p_2 = p_3 = 0$ gesetzt. Es egeben sich folgende
Bits (mit der zugewiesenen Bezeichnung darunter):

\placetable[here,none][]{}{
\starttable[|m|m|m|m|m|m|m|]
\NC 1   \NC 0   \NC 1   \NC 1   \NC 1   \NC 0   \NC 0 \SR
\NC d_1 \NC d_2 \NC d_3 \NC d_4 \NC p_1 \NC p_2 \NC p_3 \SR
\stoptable
}

Nun kann man feststellen, {\bf ob} und {\bf welches} Bit sich verändert hat.
Sollte sich der Bit $d_3$ verändert haben, werden sowohl $p_1$ als auch $p_3$
einen veränderung wiederspiegeln, wodurch $d_3$ identifierzt werden kann, da
keinem anderen Bit diese zwei Paritätbits zugewiesen sind. Dies wird auch nochmal in \in{Abbildung}[fig:mp:hammingcode:bits:trig] deutlich.

\startbuffer[fig:mp:hammingcode:bits:trig]
\startMPpage

u := 1cm;
% Farben definieren
color blue; blue := (0.2, 0.3, 0.7);
color green; green := (0.3, 0.6, 0.3);
color orange; orange := (0.8, 0.5, 0.2);
color yellow; yellow := (1, 0.8, 0.3);
color gray; gray := (0.6, 0.6, 0.6);


% Kreise definieren
path cone, ctwo, cthree;
cone := fullcircle scaled (3 * u) shifted (u * dir 0);
ctwo := fullcircle scaled (3 * u) shifted (u * dir 120);
cthree := fullcircle scaled (3 * u) shifted (u * dir 240);

% Kreise mit Transparenz und Farben füllen
fill cone withcolor blue withtransparency (1, 0.6);
fill ctwo withcolor gray withtransparency (1, 0.3);
fill cthree withcolor green withtransparency (1, 0.6);

draw cone withcolor blue;
draw ctwo withcolor gray;
draw cthree withcolor green;

label("$p_1$", (u * 1.5 * dir 0));
label("$p_2$", (u * 1.5 * dir 120));
label("$p_3$", (u * 1.5 * dir 240));

label("$1$", (u * dir 60));
label("$0$", (u * dir 180));
label("{\bfa $1$}", (u * dir 300));
label("$1$", origin);

\stopMPpage
\stopbuffer

\margintext{
\placefigure[here][fig:mp:hammingcode:bits:trig]{Verteilung der Bits in der Besipielrechnung}{
\typesetbuffer[fig:mp:hammingcode:bits:trig]
}}

\stopsectionlevel


\stopchapter

\startmode[single]

\stopbodymatter

\startbackmatter
\component c_definitions
% \component c_bibliography
\stopbackmatter

\stoptext
\stopmode
